\section{Introduction}

Polarization is a central feature of contemporary public discourse.
It appears in electoral competition, climate-policy debates, and
attitudes toward institutions, where populations are often described
as divided into mutually antagonistic camps. To study such phenomena,
one first needs a measure. Claims about the causes of polarization,
its persistence, and its possible reduction depend on what the measure
is taken to express. The problem is therefore not merely to assign a
number to a distribution of opinions, but to say what that number
means.

Several answers have been proposed. The Esteban--Ray measure
(\ERLabel)~\cite{esteban1994measurement} defines polarization
through a pairwise sum combining intra-group identification and
inter-group alienation. The Earth Mover's Distance
(\EMDLabel), from optimal-transport
theory~\cite{geher2020isometric,vallender1974calculation}, measures
the minimum cost of transforming one distribution into another.
The Tastle--Wierman Dissention measure
(\TWDLabel)~\cite{tastle2007consensus} adapts
information-theoretic ideas to ordinal scales. The Van der Eijk index
(\VdELabel)~\cite{van2001measuring} treats polarization as the
complement of agreement under a unimodality criterion. The
Koudenburg--Kiers--Kashima Index
(\KILabel)~\cite{koudenburg2021new} is an empirical
regression-based measure fitted to expert judgments of polarization.
These measures
differ in motivation and form, but they share a common feature:
polarization is computed directly from the distribution, either as a
pairwise aggregate, an ordinal dispersion index, or a transport cost.
They do not, however, identify a point at which consensus could be
reached with minimum effort, nor do they provide the corresponding
optimization structure.

We propose such a structure. Our starting point is the following
interpretation: polarization is resistance to consensus. A population
is highly polarized when much effort is needed to bring its members to
a common position; it is weakly polarized when little effort is needed.
This idea leads to an optimization problem. Given an opinion
distribution, we define its polarization as the minimum effort required
to transform it into a consensus distribution. We call this quantity
the \emph{Minimum Effort to Consensus} (MEC). The minimum is taken over
all possible consensus points. A minimizer, denoted \(y^*\), is
therefore not an external parameter, but an output of the measure: it
is the consensus point that is least costly to reach.

The definition is connected to transportation theory. In the basic
case, where effort is mass times distance, MEC is the
\(1\)-Wasserstein distance from \(\pi\) to its nearest consensus
distribution, equivalently the Earth Mover's Distance
between \(\pi\) and the least costly consensus configuration. More
generally, for parameters \(\alpha,\beta\geq 1\), MEC is the minimum,
over consensus points, of a weighted \(L^\beta\) cost in which the
weights are the \(\alpha\)-power of the population masses. This
parametric form includes two familiar dispersion measures as special
cases: mean absolute deviation about a median when
\((\alpha,\beta)=(1,1)\), and variance-like dispersion when
\((\alpha,\beta)=(1,2)\). The parameters also have a natural social
interpretation. The exponent \(\alpha\) controls the effect of group
size and corresponds to the identification dimension emphasized by
Social Identity Theory~\cite{TajfelTurner1979}; the exponent \(\beta\)
controls the cost of large opinion shifts and corresponds to the
alienation dimension, related to the cognitive cost of substantial
belief revision~\cite{Festinger1957}.

The contributions of the paper are as follows.

\paragraph{Conceptual foundation.}
We define polarization as minimum effort to consensus. Formally, this
amounts to measuring the distance from an opinion distribution to the
set of consensus distributions. The definition yields both a scalar
polarization value and an optimal consensus point. It also recovers
classical dispersion measures in special cases, while allowing two
parameters corresponding to identification and alienation.

\paragraph{Theoretical results.}
We establish a geometric characterization of MEC in the transport
case. In particular, Theorem~\ref{th:mec-as-emd} shows that MEC can be
expressed as \(\Mass{\pi}/2\) minus the Earth Mover's Distance from
\(\pi\) to the extremal distribution. Thus, in this case,
polarization is proximity to maximum disagreement. The optimization
formulation also yields structural information beyond the scalar
value. The optimal point \(y^*\) identifies the consensus position that
minimizes total effort. We prove a Minority Principle
(Theorem~\ref{thm:minority-principle-general}), according to which, in
a two-group distribution, the smaller group bears the greater share of
the effort to consensus. We also introduce a Tipping Point method
(Section~\ref{sec:tipping-point}) that gives, in closed form, the
critical mass at which moving opinions toward an extreme begins to
increase rather than decrease polarization. Together, these results
reveal that polarization is not monotone in extremism, nor even in
movement away from the optimal consensus point. Displacing the whole
of a group away from the optimal consensus point raises polarization,
but transferring only a portion of the group can lower it: the detached
fragment may be too small to exert influence at its new position,
while the original group is weakened. The Tipping Point method
makes this dependence explicit and computable. Finally,
we show that MEC
satisfies the three axioms of Esteban and Ray~\cite{esteban1994measurement},
as well as a central-split monotonicity property, under natural
conditions on the parameters.

\paragraph{Empirical and computational results.}
We compare MEC with the expert-judgment benchmark of Koudenburg,
Kiers, and Kashima~\cite{koudenburg2021new}, in which sixty researchers in
polarization studies rated fifteen hypothetical opinion distributions.
The parametrization \(\MEC[2,1.15]\) obtains Kendall's
\(\tau\approx 0.89\) against the expert ranking. This matches the
strongest parametrization of \ERLabel{} and outperforms
the dispersion-based alternatives \VdELabel{} and
\TWDLabel. MEC is also computationally
tractable. For support size \(n\), the convex objective can be
minimized by bisection in \(O(n\log(1/\varepsilon))\) time to precision
\(\varepsilon\). This is asymptotically faster than the \(O(n^2)\)
cost of pairwise measures. The difference is small on short ordinal
scales, but becomes relevant for fine-grained opinion spaces and
large-scale corpus analyses such as those in
Section~\ref{sec:experiments}.

MEC and the Esteban--Ray measure are conceptually different.
\ERLabel{} measures pairwise antagonism between groups. MEC measures
the resistance of the whole distribution to being brought into
agreement. The two measures coincide in spirit on canonical cases,
such as the symmetric two-extreme distribution, but MEC adds an
optimization structure and an endogenous consensus point while
remaining compatible with the axiomatic perspective introduced by
Esteban and Ray.

The remainder of the paper is organized as follows.
Section~\ref{sec:preliminaries} fixes notation and defines opinion
distributions and polarization measures. Section~\ref{sec:mec}
introduces MEC, motivates the parameters \(\alpha\) and \(\beta\), and
discusses computation. Section~\ref{sec:structural} proves structural
properties, including convexity, the \(L^\beta\)-dispersion
characterization, symmetry, and population scaling.
Section~\ref{sec:theoretical-applications} presents the Minority
Principle and the Tipping Point method.
Section~\ref{sec:shifting-away} proves that MEC is a polarization
measure in the standard axiomatic sense.
Section~\ref{sec:er-comparison} compares MEC with the Esteban--Ray
measure and proves axiomatic compatibility.
Section~\ref{sec:experiments} reports the empirical validation.
Section~\ref{sec:conclusion} concludes and situates the contribution
in related work.
Proofs are collected in the appendices.
